%% series_5_bessel.tex -- version ampliada de series_5 (ecuacion y funciones de Bessel)
%% Agregados: graficos de J_0, J_1, J_2, de Y_0, Y_1, de la funcion Gamma, de las sumas parciales
%% de J_0, de J_{+-1/2} (cerradas) y del comportamiento asintotico; tablas de propiedades de Gamma,
%% de que segunda solucion usar, de comportamiento, de ceros y de valores; demostraciones de las
%% Propiedades 1 y 2; ejemplos nuevos (J_{1/2}, integrales, ecuacion parametrica, modos de un tambor);
%% ejercicios con respuesta.
%% Correcciones respecto del original: se define explicitamente J_p (el original citaba una
%% ecuacion "4.6.16" que no aparecia) y se aclara que Y_p con la formula del limite vale para
%% p no entero (para p entero, incluido p=0, se usa el limite).
%% Requiere en el preambulo del documento principal:
%%   \usepackage{pgfplots}  \pgfplotsset{compat=1.18}
%% Convenciones de notacion del capitulo: en la forma estandar y''+P(x)y'+Q(x)y=0
%% los coeficientes se llaman P y Q; en la forma de Frobenius
%% x^2 y''+x p(x) y'+q(x) y=0 se usan p=(x-x0)P y q=(x-x0)^2 Q (minusculas).
 
%% --- Correccion de \ref a cajas numeradas (ver sistemas_2.tex) --------------
\makeatletter
\ifdefined\bbereffixdone\else
  \def\bbereffixdone{}
  \tcbset{before upper={%
    \ifx\thetcbcounter\@empty\else
      \protected@edef\@currentlabel{\thetcbcounter}%
      \ifdefined\phantomsection\phantomsection\fi
    \fi}}
\fi
\makeatother
 
%% --- Estilo comun de los graficos de este capitulo --------------------------
\pgfplotsset{
  bbeserie/.style={width=.94\linewidth, height=6.2cm, axis lines=middle,
    xlabel=$x$, ylabel=$y$, tick label style={font=\small}, label style={font=\small},
    legend style={font=\footnotesize, draw=bbe!40, fill=white, fill opacity=.9, text opacity=1},
    grid=major, grid style={draw=bbe!12}, clip=true,
    every axis plot/.append style={thick}},
  bbeosc/.style={bbeserie, axis lines=left}
}
\providecommand{\bbezero}[2]{\addplot[bbegray!70, thin, forget plot] coordinates {(#1,0) (#2,0)};}
\tikzset{
  bbecaja/.style={draw=bbe, fill=bbe!8, rounded corners=2pt, text=black, align=center,
    font=\small\sffamily, inner sep=4pt, text width=3.6cm, minimum height=1.1cm},
  bbeflujo/.style={-{Stealth[length=5pt]}, thick}
}
 
 
Una de las ecuaciones diferenciales más importantes en matemáticas aplicadas y física matemática es la \textbf{ecuación de Bessel de orden $p$}, definida por
\begin{equation}
x^2 y'' + xy' + (x^2 - p^2)y = 0,
\label{eq:s5-bessel}
\end{equation}
donde $p$ es una constante \textit{no negativa}. Aparece, por ejemplo, al estudiar las vibraciones de una membrana circular, la conducción del calor en un cilindro o la propagación de ondas con simetría cilíndrica. En general, no es posible obtener soluciones de forma cerrada para esta ecuación. Sin embargo, dado que $x = 0$ es un punto singular regular (aquí $p(x)\equiv1$ y $q(x)=x^2-p^2$), podemos aplicar la técnica de series de Frobenius para obtener soluciones en serie. Asumiremos que $x > 0$.
 
La ecuación indicial para \eqref{eq:s5-bessel} es
\[ r(r - 1) + r - p^2 = 0 \]
con raíces
\[ r = \pm p. \]
En consecuencia, siempre que $2p$ no sea un número entero, existirán dos soluciones en serie de Frobenius linealmente independientes. Para obtener estas soluciones, sea
\begin{equation}
y(x) = x^r \sum_{n=0}^{\infty} a_n x^n
\label{eq:s5-ansatz}
\end{equation}
de modo que
\[ y'(x) = \sum_{n=0}^{\infty} (r + n)a_n x^{r+n-1}, \qquad y''(x) = \sum_{n=0}^{\infty} (r + n)(r + n - 1)a_n x^{r+n-2}. \]
Sustituir en \eqref{eq:s5-bessel} y reordenar produce
\begin{equation}
\sum_{n=0}^{\infty} [(r + n)^2 - p^2]a_n x^{r+n} + \sum_{n=2}^{\infty} a_{n-2} x^{r+n} = 0.
\label{eq:s5-suma}
\end{equation}
Cuando $n = 0$, obtenemos la ecuación indicial cuyas raíces son, como hemos visto antes, $r = \pm p$. Cuando $n = 1$, \eqref{eq:s5-suma} implica que
\begin{equation}
[(r + 1)^2 - p^2]a_1 = 0
\label{eq:s5-n1}
\end{equation}
y para $n \ge 2$, obtenemos la relación de recurrencia general
\begin{equation}
[(r + n)^2 - p^2]a_n = -a_{n-2}, \qquad n = 2, 3, \dots
\label{eq:s5-rec}
\end{equation}
Considere la raíz $r = p$. En este caso, \eqref{eq:s5-n1} se reduce a
\[ (2p + 1)a_1 = 0 \]
de modo que, dado que $p \ge 0$,
\begin{equation}
a_1 = 0.
\label{eq:s5-a1}
\end{equation}
Establecer $r = p$ en \eqref{eq:s5-rec} produce (nótese que $(p+n)^2-p^2=n(2p+n)$)
\begin{equation}
a_n = -\frac{1}{n(2p + n)}a_{n-2}, \qquad n = 2, 3, \dots
\label{eq:s5-rec-p}
\end{equation}
Se deduce de \eqref{eq:s5-a1} y \eqref{eq:s5-rec-p} que todos los coeficientes impares son cero; es decir,
\begin{equation}
a_{2k+1} = 0, \qquad k = 0, 1, 2, \dots
\label{eq:s5-impares}
\end{equation}
Ahora considere los coeficientes pares. De \eqref{eq:s5-rec-p}, obtenemos
\[ a_2 = -\frac{1}{2(2p + 2)}a_0, \qquad a_4 = -\frac{1}{4(2p + 4)}a_2 = \frac{1}{2 \cdot 4(2p + 4)(2p + 2)}a_0, \]
y así sucesivamente. El coeficiente par general es
\[ a_{2k} = \frac{(-1)^k}{2 \cdot 4 \cdots (2k)(2p + 2)(2p + 4) \cdots (2p + 2k)}a_0, \qquad k = 1, 2, \dots, \]
que se puede escribir como
\[ a_{2k} = \frac{(-1)^k}{2^{2k}k!(p + 1)(p + 2) \cdots (p + k)}a_0, \qquad k = 1, 2, 3, \dots \]
\begin{remark}
En consecuencia, la solución en serie de Frobenius correspondiente a la ecuación de Bessel es
\begin{equation}
y_1(x) = a_0 x^p \left[ 1 + \sum_{k=1}^{\infty} \frac{(-1)^k}{2^{2k}k!(p + 1)(p + 2) \cdots (p + k)}x^{2k} \right].
\label{eq:s5-y1}
\end{equation}
Esta solución es válida para todo $x > 0$ (es más: $p,q$ son polinomios, luego la serie converge para todo $x$).
\end{remark}
 
\subsection{Funciones de Bessel}
 
Para estudiar las soluciones de la ecuación de Bessel obtenidas anteriormente, es conveniente primero reescribir \eqref{eq:s5-y1} en una forma diferente, pero equivalente. El análisis se divide en dos casos:
 
\vspace{0.3cm}
\noindent \textbf{Caso 1: $p = N$, un entero no negativo:} \quad En este caso, la solución \eqref{eq:s5-y1} se puede escribir como
\begin{equation}
y_1(x) = a_0 x^N \left[ 1 + \sum_{k=1}^{\infty} \frac{(-1)^k}{2^{2k}k!(N + 1)(N + 2) \cdots (N + k)}x^{2k} \right],
\label{eq:s5-y1-N}
\end{equation}
donde la constante $a_0$ se puede elegir arbitrariamente. Es conveniente hacer la elección
\[ a_0 = \frac{1}{N!2^N}. \]
 
\begin{definition}[Función de Bessel de primera especie, orden entero]
La solución correspondiente de la ecuación de Bessel se denota $J_N(x)$ y se llama la \textbf{función de Bessel de primera especie de orden entero $N$}. Sustituyendo $a_0$ en \eqref{eq:s5-y1-N} (y usando $N!\,(N+1)\cdots(N+k)=(N+k)!$), obtenemos
\begin{equation}
J_N(x) = \sum_{k=0}^{\infty} \frac{(-1)^k}{k!(N + k)!} \left(\frac{x}{2}\right)^{2k+N}.
\label{eq:s5-JN}
\end{equation}
\end{definition}
 
\begin{remark}
    Las funciones de Bessel de orden entero más importantes son $J_0(x)$ y $J_1(x)$, ya que, como veremos en breve, todas las demás funciones de Bessel de primera especie de orden entero se pueden expresar en términos de estas dos. Escribir los primeros términos en \eqref{eq:s5-JN} cuando $N = 0, 1$ produce, respectivamente,
\begin{align*}
J_0(x) &= 1 - \frac{1}{4}x^2 + \frac{1}{64}x^4 - \dots \\
J_1(x) &= \frac{1}{2}x \left( 1 - \frac{1}{8}x^2 + \frac{1}{192}x^4 - \dots \right).
\end{align*}
\end{remark}
 
Un análisis de estas funciones muestra que ambas oscilan con amplitud decreciente (Figura~\ref{fig:s5-J}). Además, cada una tiene un número infinito de ceros no negativos. Nótese el parecido con $\cos x$ y $\sin x$, pero con la amplitud que decae como $1/\sqrt x$ y los ceros que no son equiespaciados exactamente (más adelante veremos que sí lo son asintóticamente).
 
\begin{figure}[H]
\centering
\begin{tikzpicture}
\begin{axis}[bbeosc, width=.96\linewidth, height=6.4cm, xmin=0, xmax=15, ymin=-0.5, ymax=1.1, legend pos=north east, legend columns=3]
\bbezero{0}{15}
\addplot[bbeblue, very thick, smooth] coordinates {(0,1.0000) (0.1,0.9975) (0.2,0.9900) (0.3,0.9776) (0.4,0.9604) (0.5,0.9385) (0.6,0.9120) (0.7,0.8812) (0.8,0.8463) (0.9,0.8075) (1,0.7652) (1.1,0.7196) (1.2,0.6711) (1.3,0.6201) (1.4,0.5669) (1.5,0.5118) (1.6,0.4554) (1.7,0.3980) (1.8,0.3400) (1.9,0.2818) (2,0.2239) (2.1,0.1666) (2.2,0.1104) (2.3,0.0555) (2.4,0.0025) (2.5,-0.0484) (2.6,-0.0968) (2.7,-0.1424) (2.8,-0.1850) (2.9,-0.2243) (3,-0.2601) (3.1,-0.2921) (3.2,-0.3202) (3.3,-0.3443) (3.4,-0.3643) (3.5,-0.3801) (3.6,-0.3918) (3.7,-0.3992) (3.8,-0.4026) (3.9,-0.4018) (4,-0.3971) (4.1,-0.3887) (4.2,-0.3766) (4.3,-0.3610) (4.4,-0.3423) (4.5,-0.3205) (4.6,-0.2961) (4.7,-0.2693) (4.8,-0.2404) (4.9,-0.2097) (5,-0.1776) (5.1,-0.1443) (5.2,-0.1103) (5.3,-0.0758) (5.4,-0.0412) (5.5,-0.0068) (5.6,0.0270) (5.7,0.0599) (5.8,0.0917) (5.9,0.1220) (6,0.1506) (6.1,0.1773) (6.2,0.2017) (6.3,0.2238) (6.4,0.2433) (6.5,0.2601) (6.6,0.2740) (6.7,0.2851) (6.8,0.2931) (6.9,0.2981) (7,0.3001) (7.1,0.2991) (7.2,0.2951) (7.3,0.2882) (7.4,0.2786) (7.5,0.2663) (7.6,0.2516) (7.7,0.2346) (7.8,0.2154) (7.9,0.1944) (8,0.1717) (8.1,0.1475) (8.2,0.1222) (8.3,0.0960) (8.4,0.0692) (8.5,0.0419) (8.6,0.0146) (8.7,-0.0125) (8.8,-0.0392) (8.9,-0.0653) (9,-0.0903) (9.1,-0.1142) (9.2,-0.1367) (9.3,-0.1577) (9.4,-0.1768) (9.5,-0.1939) (9.6,-0.2090) (9.7,-0.2218) (9.8,-0.2323) (9.9,-0.2403) (10,-0.2459) (10.1,-0.2490) (10.2,-0.2496) (10.3,-0.2477) (10.4,-0.2434) (10.5,-0.2366) (10.6,-0.2276) (10.7,-0.2164) (10.8,-0.2032) (10.9,-0.1881) (11,-0.1712) (11.1,-0.1528) (11.2,-0.1330) (11.3,-0.1121) (11.4,-0.0902) (11.5,-0.0677) (11.6,-0.0446) (11.7,-0.0213) (11.8,0.0020) (11.9,0.0250) (12,0.0477) (12.1,0.0697) (12.2,0.0908) (12.3,0.1108) (12.4,0.1296) (12.5,0.1469) (12.6,0.1626) (12.7,0.1766) (12.8,0.1887) (12.9,0.1988) (13,0.2069) (13.1,0.2129) (13.2,0.2167) (13.3,0.2183) (13.4,0.2177) (13.5,0.2150) (13.6,0.2101) (13.7,0.2032) (13.8,0.1943) (13.9,0.1836) (14,0.1711) (14.1,0.1570) (14.2,0.1414) (14.3,0.1245) (14.4,0.1065) (14.5,0.0875) (14.6,0.0679) (14.7,0.0476) (14.8,0.0271) (14.9,0.0064) (15,-0.0142)};
\addplot[bbered, very thick, smooth] coordinates {(0,0.0000) (0.1,0.0499) (0.2,0.0995) (0.3,0.1483) (0.4,0.1960) (0.5,0.2423) (0.6,0.2867) (0.7,0.3290) (0.8,0.3688) (0.9,0.4059) (1,0.4401) (1.1,0.4709) (1.2,0.4983) (1.3,0.5220) (1.4,0.5419) (1.5,0.5579) (1.6,0.5699) (1.7,0.5778) (1.8,0.5815) (1.9,0.5812) (2,0.5767) (2.1,0.5683) (2.2,0.5560) (2.3,0.5399) (2.4,0.5202) (2.5,0.4971) (2.6,0.4708) (2.7,0.4416) (2.8,0.4097) (2.9,0.3754) (3,0.3391) (3.1,0.3009) (3.2,0.2613) (3.3,0.2207) (3.4,0.1792) (3.5,0.1374) (3.6,0.0955) (3.7,0.0538) (3.8,0.0128) (3.9,-0.0272) (4,-0.0660) (4.1,-0.1033) (4.2,-0.1386) (4.3,-0.1719) (4.4,-0.2028) (4.5,-0.2311) (4.6,-0.2566) (4.7,-0.2791) (4.8,-0.2985) (4.9,-0.3147) (5,-0.3276) (5.1,-0.3371) (5.2,-0.3432) (5.3,-0.3460) (5.4,-0.3453) (5.5,-0.3414) (5.6,-0.3343) (5.7,-0.3241) (5.8,-0.3110) (5.9,-0.2951) (6,-0.2767) (6.1,-0.2559) (6.2,-0.2329) (6.3,-0.2081) (6.4,-0.1816) (6.5,-0.1538) (6.6,-0.1250) (6.7,-0.0953) (6.8,-0.0652) (6.9,-0.0349) (7,-0.0047) (7.1,0.0252) (7.2,0.0543) (7.3,0.0826) (7.4,0.1096) (7.5,0.1352) (7.6,0.1592) (7.7,0.1813) (7.8,0.2014) (7.9,0.2192) (8,0.2346) (8.1,0.2476) (8.2,0.2580) (8.3,0.2657) (8.4,0.2708) (8.5,0.2731) (8.6,0.2728) (8.7,0.2697) (8.8,0.2641) (8.9,0.2559) (9,0.2453) (9.1,0.2324) (9.2,0.2174) (9.3,0.2004) (9.4,0.1816) (9.5,0.1613) (9.6,0.1395) (9.7,0.1166) (9.8,0.0928) (9.9,0.0684) (10,0.0435) (10.1,0.0184) (10.2,-0.0066) (10.3,-0.0313) (10.4,-0.0555) (10.5,-0.0789) (10.6,-0.1012) (10.7,-0.1224) (10.8,-0.1422) (10.9,-0.1603) (11,-0.1768) (11.1,-0.1913) (11.2,-0.2039) (11.3,-0.2143) (11.4,-0.2225) (11.5,-0.2284) (11.6,-0.2320) (11.7,-0.2333) (11.8,-0.2323) (11.9,-0.2290) (12,-0.2234) (12.1,-0.2157) (12.2,-0.2060) (12.3,-0.1943) (12.4,-0.1807) (12.5,-0.1655) (12.6,-0.1487) (12.7,-0.1307) (12.8,-0.1114) (12.9,-0.0912) (13,-0.0703) (13.1,-0.0489) (13.2,-0.0271) (13.3,-0.0052) (13.4,0.0166) (13.5,0.0380) (13.6,0.0590) (13.7,0.0791) (13.8,0.0984) (13.9,0.1165) (14,0.1334) (14.1,0.1488) (14.2,0.1626) (14.3,0.1747) (14.4,0.1850) (14.5,0.1934) (14.6,0.1999) (14.7,0.2043) (14.8,0.2066) (14.9,0.2069) (15,0.2051)};
\addplot[bbegreen, very thick, smooth] coordinates {(0,0.0000) (0.1,0.0012) (0.2,0.0050) (0.3,0.0112) (0.4,0.0197) (0.5,0.0306) (0.6,0.0437) (0.7,0.0588) (0.8,0.0758) (0.9,0.0946) (1,0.1149) (1.1,0.1366) (1.2,0.1593) (1.3,0.1830) (1.4,0.2074) (1.5,0.2321) (1.6,0.2570) (1.7,0.2817) (1.8,0.3061) (1.9,0.3299) (2,0.3528) (2.1,0.3746) (2.2,0.3951) (2.3,0.4139) (2.4,0.4310) (2.5,0.4461) (2.6,0.4590) (2.7,0.4696) (2.8,0.4777) (2.9,0.4832) (3,0.4861) (3.1,0.4862) (3.2,0.4835) (3.3,0.4780) (3.4,0.4697) (3.5,0.4586) (3.6,0.4448) (3.7,0.4283) (3.8,0.4093) (3.9,0.3879) (4,0.3641) (4.1,0.3383) (4.2,0.3105) (4.3,0.2811) (4.4,0.2501) (4.5,0.2178) (4.6,0.1846) (4.7,0.1506) (4.8,0.1161) (4.9,0.0813) (5,0.0466) (5.1,0.0121) (5.2,-0.0217) (5.3,-0.0547) (5.4,-0.0867) (5.5,-0.1173) (5.6,-0.1464) (5.7,-0.1737) (5.8,-0.1990) (5.9,-0.2221) (6,-0.2429) (6.1,-0.2612) (6.2,-0.2769) (6.3,-0.2899) (6.4,-0.3001) (6.5,-0.3074) (6.6,-0.3119) (6.7,-0.3135) (6.8,-0.3123) (6.9,-0.3082) (7,-0.3014) (7.1,-0.2920) (7.2,-0.2800) (7.3,-0.2656) (7.4,-0.2490) (7.5,-0.2303) (7.6,-0.2097) (7.7,-0.1875) (7.8,-0.1638) (7.9,-0.1389) (8,-0.1130) (8.1,-0.0864) (8.2,-0.0593) (8.3,-0.0320) (8.4,-0.0047) (8.5,0.0223) (8.6,0.0488) (8.7,0.0745) (8.8,0.0993) (8.9,0.1228) (9,0.1448) (9.1,0.1653) (9.2,0.1840) (9.3,0.2008) (9.4,0.2154) (9.5,0.2279) (9.6,0.2380) (9.7,0.2458) (9.8,0.2512) (9.9,0.2542) (10,0.2546) (10.1,0.2527) (10.2,0.2483) (10.3,0.2416) (10.4,0.2327) (10.5,0.2216) (10.6,0.2085) (10.7,0.1936) (10.8,0.1769) (10.9,0.1586) (11,0.1390) (11.1,0.1183) (11.2,0.0966) (11.3,0.0741) (11.4,0.0512) (11.5,0.0279) (11.6,0.0046) (11.7,-0.0185) (11.8,-0.0413) (11.9,-0.0635) (12,-0.0849) (12.1,-0.1053) (12.2,-0.1245) (12.3,-0.1424) (12.4,-0.1587) (12.5,-0.1734) (12.6,-0.1862) (12.7,-0.1972) (12.8,-0.2061) (12.9,-0.2130) (13,-0.2177) (13.1,-0.2203) (13.2,-0.2208) (13.3,-0.2191) (13.4,-0.2152) (13.5,-0.2094) (13.6,-0.2015) (13.7,-0.1917) (13.8,-0.1801) (13.9,-0.1668) (14,-0.1520) (14.1,-0.1358) (14.2,-0.1185) (14.3,-0.1001) (14.4,-0.0808) (14.5,-0.0609) (14.6,-0.0405) (14.7,-0.0199) (14.8,0.0008) (14.9,0.0214) (15,0.0416)};
\legend{$J_0$, $J_1$, $J_2$}
\end{axis}
\end{tikzpicture}
\caption{Funciones de Bessel de primera especie $J_0$, $J_1$ y $J_2$: $J_0(0)=1$, mientras que $J_n(0)=0$ para $n\ge1$ (y $J_n(x)\sim x^n$ cerca del origen).}
\label{fig:s5-J}
\end{figure}
 
La serie \eqref{eq:s5-JN} converge para todo $x$, pero cuando $x$ es grande hay que sumar muchos términos antes de que las sumas parciales se parezcan a $J_0$ (Figura~\ref{fig:s5-parciales}); por eso en la práctica se usan tablas o fórmulas asintóticas.
 
\begin{figure}[H]
\centering
\begin{tikzpicture}
\begin{axis}[bbeosc, width=.96\linewidth, height=6.4cm, xmin=0, xmax=10, ymin=-2, ymax=2, restrict y to domain=-10:10,
  legend style={at={(0.5,-0.12)}, anchor=north}, legend columns=4, samples=200]
\bbezero{0}{10}
\addplot[black, very thick, smooth] coordinates {(0,1.0000) (0.1,0.9975) (0.2,0.9900) (0.3,0.9776) (0.4,0.9604) (0.5,0.9385) (0.6,0.9120) (0.7,0.8812) (0.8,0.8463) (0.9,0.8075) (1,0.7652) (1.1,0.7196) (1.2,0.6711) (1.3,0.6201) (1.4,0.5669) (1.5,0.5118) (1.6,0.4554) (1.7,0.3980) (1.8,0.3400) (1.9,0.2818) (2,0.2239) (2.1,0.1666) (2.2,0.1104) (2.3,0.0555) (2.4,0.0025) (2.5,-0.0484) (2.6,-0.0968) (2.7,-0.1424) (2.8,-0.1850) (2.9,-0.2243) (3,-0.2601) (3.1,-0.2921) (3.2,-0.3202) (3.3,-0.3443) (3.4,-0.3643) (3.5,-0.3801) (3.6,-0.3918) (3.7,-0.3992) (3.8,-0.4026) (3.9,-0.4018) (4,-0.3971) (4.1,-0.3887) (4.2,-0.3766) (4.3,-0.3610) (4.4,-0.3423) (4.5,-0.3205) (4.6,-0.2961) (4.7,-0.2693) (4.8,-0.2404) (4.9,-0.2097) (5,-0.1776) (5.1,-0.1443) (5.2,-0.1103) (5.3,-0.0758) (5.4,-0.0412) (5.5,-0.0068) (5.6,0.0270) (5.7,0.0599) (5.8,0.0917) (5.9,0.1220) (6,0.1506) (6.1,0.1773) (6.2,0.2017) (6.3,0.2238) (6.4,0.2433) (6.5,0.2601) (6.6,0.2740) (6.7,0.2851) (6.8,0.2931) (6.9,0.2981) (7,0.3001) (7.1,0.2991) (7.2,0.2951) (7.3,0.2882) (7.4,0.2786) (7.5,0.2663) (7.6,0.2516) (7.7,0.2346) (7.8,0.2154) (7.9,0.1944) (8,0.1717) (8.1,0.1475) (8.2,0.1222) (8.3,0.0960) (8.4,0.0692) (8.5,0.0419) (8.6,0.0146) (8.7,-0.0125) (8.8,-0.0392) (8.9,-0.0653) (9,-0.0903) (9.1,-0.1142) (9.2,-0.1367) (9.3,-0.1577) (9.4,-0.1768) (9.5,-0.1939) (9.6,-0.2090) (9.7,-0.2218) (9.8,-0.2323) (9.9,-0.2403) (10,-0.2459) (10.1,-0.2490) (10.2,-0.2496) (10.3,-0.2477) (10.4,-0.2434) (10.5,-0.2366) (10.6,-0.2276) (10.7,-0.2164) (10.8,-0.2032) (10.9,-0.1881) (11,-0.1712) (11.1,-0.1528) (11.2,-0.1330) (11.3,-0.1121) (11.4,-0.0902) (11.5,-0.0677) (11.6,-0.0446) (11.7,-0.0213) (11.8,0.0020) (11.9,0.0250) (12,0.0477) (12.1,0.0697) (12.2,0.0908) (12.3,0.1108) (12.4,0.1296) (12.5,0.1469) (12.6,0.1626) (12.7,0.1766) (12.8,0.1887) (12.9,0.1988) (13,0.2069) (13.1,0.2129) (13.2,0.2167) (13.3,0.2183) (13.4,0.2177) (13.5,0.2150) (13.6,0.2101) (13.7,0.2032) (13.8,0.1943) (13.9,0.1836) (14,0.1711) (14.1,0.1570) (14.2,0.1414) (14.3,0.1245) (14.4,0.1065) (14.5,0.0875) (14.6,0.0679) (14.7,0.0476) (14.8,0.0271) (14.9,0.0064) (15,-0.0142)};
\addplot[bbered, domain=0:10] {1 - (x/2)^2 + (x/2)^4/4};
\addplot[bbeamber, domain=0:10] {1 - (x/2)^2 + (x/2)^4/4 - (x/2)^6/36 + (x/2)^8/576};
\addplot[bbegreen, domain=0:10] {1 - (x/2)^2 + (x/2)^4/4 - (x/2)^6/36 + (x/2)^8/576 - (x/2)^10/14400 + (x/2)^12/518400};
\legend{$J_0$, hasta $x^4$, hasta $x^8$, hasta $x^{12}$}
\end{axis}
\end{tikzpicture}
\caption{Sumas parciales de la serie de $J_0$. Con más términos el intervalo de buen ajuste crece, pero la serie no es práctica para $x$ grande.}
\label{fig:s5-parciales}
\end{figure}
 
\vspace{0.3cm}
\noindent \textbf{Caso 2: $p$ un número no entero:} \quad Para obtener una fórmula análoga a \eqref{eq:s5-JN} para la solución en serie de Frobenius \eqref{eq:s5-y1} cuando $p$ no es un entero, necesitamos introducir la \textit{función gamma}. Esta función se puede considerar como la generalización de la función factorial al caso de los números reales no enteros.
 
\begin{definition}[Función gamma]
    La función gamma, $\Gamma$, se define por
\[ \Gamma(p) = \int_0^{\infty} t^{p-1} e^{-t} \, dt, \qquad p > 0. \]
\end{definition}
 
La Tabla~\ref{tab:s5-gamma} reúne las propiedades que usaremos; la Figura~\ref{fig:s5-gamma} muestra cómo $\Gamma$ interpola los factoriales. La propiedad clave es $\Gamma(p+1)=p\,\Gamma(p)$ (se obtiene integrando por partes), que permite escribir $\Gamma(p+k+1)=(p+k)\cdots(p+1)\Gamma(p+1)$.
 
\begin{table}[H]
\centering
\small
\caption{Propiedades de la función gamma.}
\label{tab:s5-gamma}
\begin{tabular}{|l|l|}
\hline
\rowcolor{bbe!15}\textbf{Propiedad} & \textbf{Valores} \\ \hline
Ecuación funcional & $\Gamma(p+1)=p\,\Gamma(p)$ \\ \hline
Enteros & $\Gamma(1)=1$, \quad $\Gamma(n+1)=n!$ \\ \hline
Semienteros & $\Gamma(\tfrac12)=\sqrt\pi$, \quad $\Gamma(\tfrac32)=\tfrac{\sqrt\pi}{2}$, \quad $\Gamma(\tfrac52)=\tfrac{3\sqrt\pi}{4}$ \\ \hline
Extensión a $p<0$ (no entero) & $\Gamma(p)=\dfrac{\Gamma(p+1)}{p}$; en $p=0,-1,-2,\dots$ hay polos, es decir, $\dfrac{1}{\Gamma(p)}=0$ \\ \hline
\end{tabular}
\end{table}
 
\begin{figure}[H]
\centering
\begin{tikzpicture}
\begin{axis}[bbeserie, width=.8\linewidth, height=5.6cm, xmin=0, xmax=5, ymin=0, ymax=8, restrict y to domain=0:30, xtick={1,2,3,4,5}]
\addplot[bbeblue, very thick, smooth] coordinates {(0.2,4.5908) (0.25,3.6256) (0.3,2.9916) (0.35,2.5461) (0.4,2.2182) (0.45,1.9681) (0.5,1.7725) (0.55,1.6161) (0.6,1.4892) (0.65,1.3848) (0.7,1.2981) (0.75,1.2254) (0.8,1.1642) (0.85,1.1125) (0.9,1.0686) (0.95,1.0315) (1,1.0000) (1.05,0.9735) (1.1,0.9514) (1.15,0.9330) (1.2,0.9182) (1.25,0.9064) (1.3,0.8975) (1.35,0.8912) (1.4,0.8873) (1.45,0.8857) (1.5,0.8862) (1.55,0.8889) (1.6,0.8935) (1.65,0.9001) (1.7,0.9086) (1.75,0.9191) (1.8,0.9314) (1.85,0.9456) (1.9,0.9618) (1.95,0.9799) (2,1.0000) (2.05,1.0222) (2.1,1.0465) (2.15,1.0730) (2.2,1.1018) (2.25,1.1330) (2.3,1.1667) (2.35,1.2031) (2.4,1.2422) (2.45,1.2842) (2.5,1.3293) (2.55,1.3777) (2.6,1.4296) (2.65,1.4852) (2.7,1.5447) (2.75,1.6084) (2.8,1.6765) (2.85,1.7494) (2.9,1.8274) (2.95,1.9108) (3,2.0000) (3.05,2.0955) (3.1,2.1976) (3.15,2.3069) (3.2,2.4240) (3.25,2.5493) (3.3,2.6834) (3.35,2.8272) (3.4,2.9812) (3.45,3.1463) (3.5,3.3234) (3.55,3.5133) (3.6,3.7170) (3.65,3.9358) (3.7,4.1707) (3.75,4.4230) (3.8,4.6942) (3.85,4.9857) (3.9,5.2993) (3.95,5.6368) (4,6.0000) (4.05,6.3912) (4.1,6.8126) (4.15,7.2669) (4.2,7.7567) (4.25,8.2851) (4.3,8.8553) (4.35,9.4710) (4.4,10.1361) (4.45,10.8548) (4.5,11.6317) (4.55,12.4720) (4.6,13.3813) (4.65,14.3655) (4.7,15.4314) (4.75,16.5862) (4.8,17.8379) (4.85,19.1951) (4.9,20.6674) (4.95,22.2652) (5,24.0000)};
\addplot[only marks, mark=*, mark size=2.2pt, bbered] coordinates {(1,1) (2,1) (3,2) (4,6) (0.5,1.7725) (1.5,0.8862) (2.5,1.3293)};
\end{axis}
\end{tikzpicture}
\caption{$\Gamma(p)$ en $p>0$. Los puntos rojos son $\Gamma(1)=1$, $\Gamma(2)=1$, $\Gamma(3)=2$, $\Gamma(4)=6$ y los valores en $\tfrac12,\tfrac32,\tfrac52$.}
\label{fig:s5-gamma}
\end{figure}
 
Ahora, para $p\ge0$ cualquiera, elegimos $a_0=\dfrac{1}{2^p\,\Gamma(p+1)}$ en \eqref{eq:s5-y1}. Como $a_{2k}x^{2k+p}=\dfrac{(-1)^k}{2^{2k+p}\,k!\,\Gamma(p+k+1)}x^{2k+p}$, obtenemos:
 
\begin{definition}[Función de Bessel de primera especie]
Para $p\ge0$,
\begin{equation}
J_p(x) = \sum_{k=0}^{\infty} \frac{(-1)^k}{\Gamma(k + 1)\Gamma(p + k + 1)} \left( \frac{x}{2} \right)^{2k+p}.
\label{eq:s5-Jp}
\end{equation}
Cuando $p=N$ es entero, $\Gamma(N+k+1)=(N+k)!$ y se recupera \eqref{eq:s5-JN}.
\end{definition}
 
Ahora consideremos determinar la solución general de la ecuación de Bessel. Para todo $p \ge 0$, hemos demostrado que una solución a la ecuación de Bessel en $(0, \infty)$ está dada en \eqref{eq:s5-Jp}. Por lo tanto, requerimos una segunda solución linealmente independiente. Dado que las raíces de la ecuación indicial son $r = \pm p$, se deduce de nuestra teoría general de Frobenius que, siempre que $2p$ no sea igual a un número entero, existirá una segunda solución en serie de Frobenius linealmente independiente correspondiente a la raíz $r = -p$. No es demasiado difícil ver que esta solución se puede obtener reemplazando $p$ por $-p$ en \eqref{eq:s5-Jp}. Así, obtenemos
\begin{equation}
J_{-p}(x) = \sum_{k=0}^{\infty} \frac{(-1)^k}{\Gamma(k + 1)\Gamma(k - p + 1)} \left(\frac{x}{2}\right)^{2k-p}.
\label{eq:s5-Jmp}
\end{equation}
 
\begin{remark}
En consecuencia, la solución general a la ecuación de Bessel de orden $p$ cuando $2p$ no es igual a un entero es
\begin{equation}
y(x) = c_1 J_p(x) + c_2 J_{-p}(x),
\label{eq:s5-gen-Jmp}
\end{equation}
válida para todo $x>0$.
\end{remark}
 
Cuando $2p$ es igual a un entero, surgen dos subcasos dependiendo de si $p$ es en sí mismo un entero o un semi-entero (es decir, $\frac{1}{2}, \frac{3}{2}, \dots$). En el último caso, con $p$ de la forma $(2j + 1)/2$ ($j$ un entero no negativo), un análisis directo de la relación de recurrencia \eqref{eq:s5-rec} con $r = -p$ muestra que también existe una segunda serie de Frobenius en este caso (es el Caso III con $\alpha=0$, como en el Ejemplo~\ref{ex:s4-sin-log}), y está, de hecho, dada por \eqref{eq:s5-Jmp}. Así, \eqref{eq:s5-gen-Jmp} representa la solución general de la ecuación de Bessel, siempre que $p$ no sea igual a un número entero.
 
\begin{example}\label{ex:s5-medio}
\textbf{(Orden $\frac12$: las funciones de Bessel se vuelven elementales.)} Muestre que $J_{1/2}(x)=\sqrt{\dfrac{2}{\pi x}}\,\sin x$ y $J_{-1/2}(x)=\sqrt{\dfrac{2}{\pi x}}\,\cos x$.
 
\noindent \textbf{Solución:} \quad De \eqref{eq:s5-Jp} con $p=\frac12$ y la propiedad $\Gamma(\tfrac12)=\sqrt\pi$ se obtiene $\Gamma(k+\tfrac32)=\dfrac{(2k+1)!\,\sqrt\pi}{2\cdot4^k\,k!}$, de modo que $k!\,\Gamma(k+\tfrac32)=\dfrac{(2k+1)!\sqrt\pi}{2\cdot4^k}$. Entonces
\[ J_{1/2}(x)=\sum_{k=0}^{\infty}\frac{(-1)^k\,2\cdot4^k}{(2k+1)!\sqrt\pi}\Bigl(\frac x2\Bigr)^{2k+\frac12}
=\sqrt{\frac{2}{\pi x}}\sum_{k=0}^{\infty}\frac{(-1)^k x^{2k+1}}{(2k+1)!}=\sqrt{\frac{2}{\pi x}}\,\sin x. \]
El caso $J_{-1/2}$ es análogo y produce la serie del coseno. \hfill $\square$
\end{example}
 
\begin{figure}[H]
\centering
\begin{tikzpicture}
\begin{axis}[bbeosc, width=.96\linewidth, height=6cm, xmin=0, xmax=14, ymin=-1.2, ymax=1.8, restrict y to domain=-4:4,
  legend pos=north east, legend columns=3, samples=300]
\bbezero{0}{14}
\addplot[bbeblue, very thick, domain=0.05:14] {sqrt(2/(pi*x))*sin(deg(x))};
\addplot[bbered, very thick, domain=0.05:14] {sqrt(2/(pi*x))*cos(deg(x))};
\addplot[bbegray, dashed, domain=0.05:14] {sqrt(2/(pi*x))};
\addplot[bbegray, dashed, domain=0.05:14, forget plot] {-sqrt(2/(pi*x))};
\legend{$J_{1/2}$, $J_{-1/2}$, $\pm\sqrt{2/(\pi x)}$}
\end{axis}
\end{tikzpicture}
\caption{$J_{\pm1/2}$ en términos de $\sin x$ y $\cos x$: $J_{1/2}$ está acotada en $0$ mientras que $J_{-1/2}\sim x^{-1/2}$ explota; ambas oscilan dentro de la envolvente $\pm\sqrt{2/(\pi x)}$.}
\label{fig:s5-medio}
\end{figure}
 
Para $p$ entero la situación es distinta: la fórmula \eqref{eq:s5-Jmp} no da una solución independiente. En efecto, si $p=n$ es entero, $\dfrac{1}{\Gamma(k-n+1)}=0$ para $k=0,\dots,n-1$ (polos de $\Gamma$), de modo que la serie empieza en $k=n$ y, con el cambio $k=m+n$, se obtiene
\[ J_{-n}(x)=(-1)^nJ_n(x), \]
es decir, $J_{-n}$ y $J_n$ son linealmente \emph{dependientes}. La Tabla~\ref{tab:s5-segunda} resume qué segunda solución usar.
 
En la práctica, en lugar de usar $J_{-p}$ como la segunda solución linealmente independiente de la ecuación de Bessel, es usual usar la siguiente combinación lineal de estas dos soluciones, válida para $p$ no entero:
\begin{equation}
Y_p(x) = \frac{J_p(x) \cos p\pi - J_{-p}(x)}{\sin p\pi}.
\label{eq:s5-Yp}
\end{equation}
 
La función definida en \eqref{eq:s5-Yp} se llama la \textbf{función de Bessel de segunda especie de orden $p$}. Usando $Y_p$, podemos por lo tanto escribir la solución general de la ecuación de Bessel, cuando $p$ no es un entero, en la forma
\begin{equation}
y(x) = c_1 J_p(x) + c_2 Y_p(x).
\label{eq:s5-gen-Yp}
\end{equation}
 
La determinación de una segunda solución linealmente independiente para la ecuación de Bessel cuando $p$ es un entero $n\ge0$ es un poco más complicada. De nuestra teoría general de Frobenius, ciertamente conocemos la forma de la segunda solución linealmente independiente: Caso II si $n=0$ y Caso III con $\alpha\neq0$ si $n\ge1$, a saber,
\[ y_2(x) = A J_n(x) \ln x + x^{-n} \sum_{k=0}^{\infty} b_k x^k \]
(para $n\ge1$ el logaritmo no puede evitarse, pues la condición de consistencia da $\alpha\neq0$) y los coeficientes $A, b_k$ podrían determinarse por sustitución directa en \eqref{eq:s5-bessel}.
 
\begin{definition}[Función de Bessel de segunda especie, orden entero]
Extendemos la definición de la función de Bessel de segunda especie al caso de enteros $n\ge0$ mediante
\begin{equation}
Y_n(x) = \lim_{p \to n} \left[ \frac{J_p(x) \cos p\pi - J_{-p}(x)}{\sin p\pi} \right].
\label{eq:s5-Yn-lim}
\end{equation}
\end{definition}
 
Se puede demostrar que el límite anterior existe y que la función resultante es, de hecho, una segunda solución linealmente independiente de la ecuación de Bessel. Podríamos derivar la representación en serie de \eqref{eq:s5-Yn-lim} evaluando el límite explícitamente. Los cálculos son largos y tediosos, por lo que los omitimos. El resultado de estos cálculos es
\begin{align*}
Y_n(x) &= \frac{2}{\pi} J_n(x) \left[ \ln \left( \frac{x}{2} \right) + \gamma \right] - \frac{1}{\pi} \left[ \sum_{k=0}^{n-1} \frac{(n - k - 1)!}{k!} \left( \frac{x}{2} \right)^{2k-n} + \frac{s_n}{n!} \left( \frac{x}{2} \right)^n \right. \\
&\quad \left. + \sum_{k=1}^{\infty} (-1)^k \frac{s_k + s_{n+k}}{k!(n + k)!} \left( \frac{x}{2} \right)^{2k+n} \right],
\end{align*}
donde
\[ s_k = 1 + \frac{1}{2} + \frac{1}{3} + \dots + \frac{1}{k} \]
y
\[ \gamma = \lim_{k \to \infty} (s_k - \ln k) \approx 0.577215664\dots \]
El valor $\gamma$ se llama \textit{constante de Euler}. Así, \eqref{eq:s5-gen-Yp} (con $Y_n$ definida por el límite) representa la solución general de la ecuación de Bessel de orden arbitrario $p$.
 
\begin{table}[H]
\centering
\small
\caption{Qué segunda solución linealmente independiente usar en la ecuación de Bessel de orden $p\ge0$.}
\label{tab:s5-segunda}
\begin{tabular}{|p{3.6cm}|p{3.8cm}|p{3.2cm}|p{3.4cm}|}
\hline
\rowcolor{bbe!15}\textbf{Orden $p$} & \textbf{Raíces $\pm p$ (diferencia $2p$)} & \textbf{Segunda solución} & \textbf{Caso de Frobenius} \\ \hline
$2p\notin\mathbb Z$ & no entero & $J_{-p}$ (o $Y_p$) & I \\ \hline
$p=\frac12,\frac32,\dots$ & entero impar & $J_{-p}$ (o $Y_p$) & III con $\alpha=0$ (sin $\ln$) \\ \hline
$p=0$ & raíz doble & $Y_0$ & II (con $\ln$) \\ \hline
$p=1,2,\dots$ & entero par & $Y_p$; ¡$J_{-p}=(-1)^pJ_p$ no sirve! & III con $\alpha\ne0$ (con $\ln$) \\ \hline
\end{tabular}
\end{table}
 
\vspace{0.3cm}
 
En la práctica, usualmente estamos interesados en soluciones de la ecuación de Bessel que están acotadas en $x = 0$. Sin embargo, las funciones de Bessel de segunda especie siempre están no acotadas en $x = 0$ (Figura~\ref{fig:s5-Y}). (Cuando $p$ no es un entero, esto se deduce del hecho de que $x^{-p}$ tiene un exponente negativo, y cuando $p$ es un entero, se debe al segundo término en la expansión de $Y_n$ dada anteriormente). Por lo tanto, usualmente requerimos solo las funciones de Bessel de primera especie.
 
\begin{figure}[H]
\centering
\begin{tikzpicture}
\begin{axis}[bbeosc, width=.96\linewidth, height=6.2cm, xmin=0, xmax=15, ymin=-2, ymax=1, legend pos=south east, legend columns=2]
\bbezero{0}{15}
\addplot[bbeblue, very thick, smooth] coordinates {(0.1,-1.5342) (0.2,-1.0811) (0.3,-0.8073) (0.4,-0.6060) (0.5,-0.4445) (0.6,-0.3085) (0.7,-0.1907) (0.8,-0.0868) (0.9,0.0056) (1,0.0883) (1.1,0.1622) (1.2,0.2281) (1.3,0.2865) (1.4,0.3379) (1.5,0.3824) (1.6,0.4204) (1.7,0.4520) (1.8,0.4774) (1.9,0.4968) (2,0.5104) (2.1,0.5183) (2.2,0.5208) (2.3,0.5181) (2.4,0.5104) (2.5,0.4981) (2.6,0.4813) (2.7,0.4605) (2.8,0.4359) (2.9,0.4079) (3,0.3769) (3.1,0.3431) (3.2,0.3071) (3.3,0.2691) (3.4,0.2296) (3.5,0.1890) (3.6,0.1477) (3.7,0.1061) (3.8,0.0645) (3.9,0.0234) (4,-0.0169) (4.1,-0.0561) (4.2,-0.0938) (4.3,-0.1296) (4.4,-0.1633) (4.5,-0.1947) (4.6,-0.2235) (4.7,-0.2494) (4.8,-0.2723) (4.9,-0.2921) (5,-0.3085) (5.1,-0.3216) (5.2,-0.3313) (5.3,-0.3374) (5.4,-0.3402) (5.5,-0.3395) (5.6,-0.3354) (5.7,-0.3282) (5.8,-0.3177) (5.9,-0.3044) (6,-0.2882) (6.1,-0.2694) (6.2,-0.2483) (6.3,-0.2251) (6.4,-0.1999) (6.5,-0.1732) (6.6,-0.1452) (6.7,-0.1162) (6.8,-0.0864) (6.9,-0.0563) (7,-0.0259) (7.1,0.0042) (7.2,0.0339) (7.3,0.0628) (7.4,0.0907) (7.5,0.1173) (7.6,0.1424) (7.7,0.1658) (7.8,0.1872) (7.9,0.2065) (8,0.2235) (8.1,0.2381) (8.2,0.2501) (8.3,0.2595) (8.4,0.2662) (8.5,0.2702) (8.6,0.2715) (8.7,0.2700) (8.8,0.2659) (8.9,0.2592) (9,0.2499) (9.1,0.2383) (9.2,0.2245) (9.3,0.2086) (9.4,0.1907) (9.5,0.1712) (9.6,0.1502) (9.7,0.1279) (9.8,0.1045) (9.9,0.0804) (10,0.0557) (10.1,0.0307) (10.2,0.0056) (10.3,-0.0193) (10.4,-0.0437) (10.5,-0.0675) (10.6,-0.0904) (10.7,-0.1122) (10.8,-0.1326) (10.9,-0.1516) (11,-0.1688) (11.1,-0.1843) (11.2,-0.1977) (11.3,-0.2091) (11.4,-0.2183) (11.5,-0.2252) (11.6,-0.2299) (11.7,-0.2322) (11.8,-0.2322) (11.9,-0.2298) (12,-0.2252) (12.1,-0.2184) (12.2,-0.2095) (12.3,-0.1986) (12.4,-0.1858) (12.5,-0.1712) (12.6,-0.1551) (12.7,-0.1375) (12.8,-0.1187) (12.9,-0.0989) (13,-0.0782) (13.1,-0.0569) (13.2,-0.0352) (13.3,-0.0134) (13.4,0.0085) (13.5,0.0301) (13.6,0.0512) (13.7,0.0717) (13.8,0.0913) (13.9,0.1099) (14,0.1272) (14.1,0.1431) (14.2,0.1575) (14.3,0.1703) (14.4,0.1812) (14.5,0.1903) (14.6,0.1974) (14.7,0.2025) (14.8,0.2056) (14.9,0.2065) (15,0.2055)};
\addplot[bbered, very thick, smooth] coordinates {(0.3,-2.2931) (0.4,-1.7809) (0.5,-1.4715) (0.6,-1.2604) (0.7,-1.1032) (0.8,-0.9781) (0.9,-0.8731) (1,-0.7812) (1.1,-0.6981) (1.2,-0.6211) (1.3,-0.5485) (1.4,-0.4791) (1.5,-0.4123) (1.6,-0.3476) (1.7,-0.2847) (1.8,-0.2237) (1.9,-0.1644) (2,-0.1070) (2.1,-0.0517) (2.2,0.0015) (2.3,0.0523) (2.4,0.1005) (2.5,0.1459) (2.6,0.1884) (2.7,0.2276) (2.8,0.2635) (2.9,0.2959) (3,0.3247) (3.1,0.3496) (3.2,0.3707) (3.3,0.3879) (3.4,0.4010) (3.5,0.4102) (3.6,0.4154) (3.7,0.4167) (3.8,0.4141) (3.9,0.4078) (4,0.3979) (4.1,0.3846) (4.2,0.3680) (4.3,0.3484) (4.4,0.3260) (4.5,0.3010) (4.6,0.2737) (4.7,0.2445) (4.8,0.2136) (4.9,0.1812) (5,0.1479) (5.1,0.1137) (5.2,0.0792) (5.3,0.0445) (5.4,0.0101) (5.5,-0.0238) (5.6,-0.0568) (5.7,-0.0887) (5.8,-0.1192) (5.9,-0.1481) (6,-0.1750) (6.1,-0.1998) (6.2,-0.2223) (6.3,-0.2422) (6.4,-0.2596) (6.5,-0.2741) (6.6,-0.2857) (6.7,-0.2945) (6.8,-0.3002) (6.9,-0.3029) (7,-0.3027) (7.1,-0.2995) (7.2,-0.2934) (7.3,-0.2846) (7.4,-0.2731) (7.5,-0.2591) (7.6,-0.2428) (7.7,-0.2243) (7.8,-0.2039) (7.9,-0.1817) (8,-0.1581) (8.1,-0.1331) (8.2,-0.1072) (8.3,-0.0806) (8.4,-0.0535) (8.5,-0.0262) (8.6,0.0011) (8.7,0.0280) (8.8,0.0544) (8.9,0.0799) (9,0.1043) (9.1,0.1275) (9.2,0.1491) (9.3,0.1691) (9.4,0.1871) (9.5,0.2032) (9.6,0.2171) (9.7,0.2287) (9.8,0.2379) (9.9,0.2447) (10,0.2490) (10.1,0.2508) (10.2,0.2502) (10.3,0.2471) (10.4,0.2416) (10.5,0.2337) (10.6,0.2236) (10.7,0.2114) (10.8,0.1973) (10.9,0.1813) (11,0.1637) (11.1,0.1446) (11.2,0.1243) (11.3,0.1029) (11.4,0.0807) (11.5,0.0579) (11.6,0.0348) (11.7,0.0114) (11.8,-0.0118) (11.9,-0.0347) (12,-0.0571) (12.1,-0.0787) (12.2,-0.0994) (12.3,-0.1189) (12.4,-0.1371) (12.5,-0.1538) (12.6,-0.1689) (12.7,-0.1821) (12.8,-0.1935) (12.9,-0.2028) (13,-0.2101) (13.1,-0.2152) (13.2,-0.2182) (13.3,-0.2190) (13.4,-0.2176) (13.5,-0.2140) (13.6,-0.2084) (13.7,-0.2007) (13.8,-0.1912) (13.9,-0.1798) (14,-0.1666) (14.1,-0.1520) (14.2,-0.1359) (14.3,-0.1186) (14.4,-0.1003) (14.5,-0.0810) (14.6,-0.0612) (14.7,-0.0408) (14.8,-0.0202) (14.9,0.0005) (15,0.0211)};
\legend{$Y_0$, $Y_1$}
\end{axis}
\end{tikzpicture}
\caption{Funciones de Bessel de segunda especie: $Y_0(x)\sim\frac2\pi\ln x$ y $Y_1(x)\sim-\frac{2}{\pi x}$ tienden a $-\infty$ cuando $x\to0^+$.}
\label{fig:s5-Y}
\end{figure}
 
\begin{table}[H]
\centering
\small
\caption{Comportamiento de $J_p$ y $Y_p$ cerca del origen y en el infinito.}
\label{tab:s5-comportamiento}
\begin{tabular}{|c|c|c|}
\hline
\rowcolor{bbe!15} & $\mathbf{x\to0^+}$ & $\mathbf{x\to\infty}$ \\ \hline
$J_p(x)$ & $\dfrac{1}{\Gamma(p+1)}\Bigl(\dfrac x2\Bigr)^p$ \quad (\emph{acotada}) & $\sqrt{\dfrac{2}{\pi x}}\cos\!\Bigl(x-\dfrac{p\pi}{2}-\dfrac\pi4\Bigr)$ \\ \hline
$Y_0(x)$ & $\dfrac2\pi\ln x$ \quad (\emph{no acotada}) & $\sqrt{\dfrac{2}{\pi x}}\sin\!\Bigl(x-\dfrac\pi4\Bigr)$ \\ \hline
$Y_p(x)$, $p>0$ & $-\dfrac{\Gamma(p)}{\pi}\Bigl(\dfrac 2x\Bigr)^p$ \quad (\emph{no acotada}) & $\sqrt{\dfrac{2}{\pi x}}\sin\!\Bigl(x-\dfrac{p\pi}{2}-\dfrac\pi4\Bigr)$ \\ \hline
\end{tabular}
\end{table}
 
En esta sección, enumeramos varias propiedades de las funciones de Bessel de primera especie que ayudan ya sea a tabular valores de las funciones o a trabajar con las funciones mismas.
 
\subsection{Propiedades de las funciones de Bessel}
 
Primero recordamos de \eqref{eq:s5-Jp} la definición de las funciones de Bessel de primera especie de orden $p$, a saber
\begin{equation}
J_p(x) = \sum_{k=0}^{\infty} \frac{(-1)^k}{\Gamma(k + 1)\Gamma(p + k + 1)} \left( \frac{x}{2} \right)^{2k+p}.
\label{eq:s5-Jp-repetida}
\end{equation}
 
\begin{theorem}[Propiedad 1]\label{thm:s5-P1}
\begin{equation}
\frac{d}{dx} [x^p J_p(x)] = x^p J_{p-1}(x).
\label{eq:s5-P1}
\end{equation}
\end{theorem}
 
\begin{theorem}[Propiedad 2]\label{thm:s5-P2}
\begin{equation}
\frac{d}{dx} [x^{-p} J_p(x)] = -x^{-p} J_{p+1}(x).
\label{eq:s5-P2}
\end{equation}
\end{theorem}
 
\begin{remark}
\textbf{Demostración de la Propiedad 1.} De \eqref{eq:s5-Jp-repetida}, $x^pJ_p(x)=\sum_{k\ge0}\dfrac{(-1)^k\,x^{2k+2p}}{2^{2k+p}\,k!\,\Gamma(k+p+1)}$. Derivando término a término aparece el factor $2(k+p)$, y $\dfrac{2(k+p)}{\Gamma(k+p+1)}=\dfrac{2}{\Gamma(k+p)}$ (ecuación funcional), de modo que
\[ \frac{d}{dx}[x^pJ_p]=\sum_{k\ge0}\frac{(-1)^k\,2\,x^{2k+2p-1}}{2^{2k+p}\,k!\,\Gamma(k+p)}=x^p\sum_{k\ge0}\frac{(-1)^k}{k!\,\Gamma(k+p)}\Bigl(\frac x2\Bigr)^{2k+p-1}=x^pJ_{p-1}(x). \]
La Propiedad 2 se demuestra igual: en $x^{-p}J_p=\sum_k\frac{(-1)^kx^{2k}}{2^{2k+p}k!\Gamma(k+p+1)}$ el término $k=0$ es constante, el factor $2k/k!=2/(k-1)!$ permite reindexar $k=j+1$ y el resultado es $-x^{-p}J_{p+1}(x)$.
\end{remark}
 
Ahora derivamos dos identidades satisfechas por las derivadas de $J_p$. Expandir las derivadas en los lados derechos de \eqref{eq:s5-P1} y \eqref{eq:s5-P2} y dividir las ecuaciones resultantes por $x^p$ y $x^{-p}$ respectivamente produce
\begin{align}
J_p'(x) + x^{-1}p J_p(x) &= J_{p-1}(x) \label{eq:s5-25}\\
J_p'(x) - x^{-1}p J_p(x) &= -J_{p+1}(x). \label{eq:s5-26}
\end{align}
Restando \eqref{eq:s5-26} de \eqref{eq:s5-25} y reordenando los términos, obtenemos:
 
\begin{theorem}[Propiedad 3: relación de recurrencia]\label{thm:s5-P3}
\begin{equation}
J_{p+1}(x) = 2x^{-1}p J_p(x) - J_{p-1}(x).
\label{eq:s5-P3}
\end{equation}
\end{theorem}
 
De manera similar, sumar \eqref{eq:s5-25} y \eqref{eq:s5-26} y reordenar produce:
 
\begin{theorem}[Propiedad 4]\label{thm:s5-P4}
\begin{equation}
J_p'(x) = \frac{1}{2} [J_{p-1}(x) - J_{p+1}(x)].
\label{eq:s5-P4}
\end{equation}
\end{theorem}
 
Estas fórmulas nos permiten expresar funciones de Bessel de orden superior y sus derivadas en términos de funciones de orden inferior. Por ejemplo, todas las funciones de Bessel de orden entero se pueden expresar en términos de $J_0(x)$ y $J_1(x)$.
 
\begin{example}\label{ex:s5-J2J3}
  Exprese $J_2(x)$ y $J_3(x)$ en términos de $J_0(x)$ y $J_1(x)$.
 
\noindent \textbf{Solución:} \quad Aplicando el Teorema~\ref{thm:s5-P3} con $p = 1$, obtenemos
\[ J_2(x) = 2x^{-1}J_1(x) - J_0(x). \]
De manera similar, cuando $p = 2$,
\[ J_3(x) = 4x^{-1}J_2(x) - J_1(x) = x^{-2}(8 - x^2)J_1(x) - 4x^{-1}J_0(x). \tag*{$\square$} \]
\end{example}
 
\begin{example}\label{ex:s5-integral}
\textbf{(Integrales y derivadas.)} (a) Calcule $\displaystyle\int_0^3 xJ_0(x)\,dx$. (b) Muestre que $J_0'(x)=-J_1(x)$. (c) Calcule $\displaystyle\int x^2J_1(x)\,dx$.
 
\noindent \textbf{Solución:} \quad (a) La Propiedad 1 con $p=1$ dice $\frac{d}{dx}[xJ_1(x)]=xJ_0(x)$, luego
\[ \int_0^3xJ_0(x)\,dx=\bigl[xJ_1(x)\bigr]_0^3=3J_1(3)\approx1{,}0172. \]
(b) La Propiedad 2 con $p=0$ dice $\frac{d}{dx}[J_0(x)]=-J_1(x)$. (c) La Propiedad 1 con $p=2$ dice $\frac{d}{dx}[x^2J_2(x)]=x^2J_1(x)$, de modo que $\int x^2J_1(x)\,dx=x^2J_2(x)+C$. \hfill $\square$
\end{example}
 
\begin{example}\label{ex:s5-parametrica}
\textbf{(Una ecuación que se reduce a Bessel.)} Resuelva $x^2y''+xy'+(9x^2-4)y=0$, $x>0$.
 
\noindent \textbf{Solución:} \quad Sea $z=3x$ y $y(x)=u(z)$. Entonces $xy'=zu'$ y $x^2y''=z^2u''$, y la ecuación queda $z^2u''+zu'+(z^2-4)u=0$: Bessel de orden $p=2$. Por lo tanto
\[ y(x)=c_1J_2(3x)+c_2Y_2(3x), \]
y la solución acotada en $x=0$ es múltiplo de $J_2(3x)$. En general, $x^2y''+xy'+(\lambda^2x^2-p^2)y=0$ tiene solución general $c_1J_p(\lambda x)+c_2Y_p(\lambda x)$. \hfill $\square$
\end{example}
 
\subsubsection*{Ceros y comportamiento para $x$ grande}
 
Como $J_p(x)\approx\sqrt{2/(\pi x)}\cos\bigl(x-\frac{p\pi}2-\frac\pi4\bigr)$ para $x$ grande, los ceros de $J_p$ quedan casi equiespaciados, con separación cercana a $\pi$. Por ejemplo, para $p=0$ los ceros son aproximadamente $(k-\tfrac14)\pi$ y para $p=1$, $(k+\tfrac14)\pi$ (Tabla~\ref{tab:s5-ceros} y Figura~\ref{fig:s5-asint}).
 
\begin{table}[H]
\centering
\small
\caption{Primeros cinco ceros positivos de $J_0$ y $J_1$, y su aproximación asintótica.}
\label{tab:s5-ceros}
\begin{tabular}{|c|c|c|c|c|}
\hline
\rowcolor{bbe!15}$k$ & $j_{0,k}$ (exacto) & $(k-\frac14)\pi$ & $j_{1,k}$ (exacto) & $(k+\frac14)\pi$ \\ \hline
$1$ & $2{,}4048$ & $2{,}3562$ & $3{,}8317$ & $3{,}9270$ \\ \hline
$2$ & $5{,}5201$ & $5{,}4978$ & $7{,}0156$ & $7{,}0686$ \\ \hline
$3$ & $8{,}6537$ & $8{,}6394$ & $10{,}1735$ & $10{,}2102$ \\ \hline
$4$ & $11{,}7915$ & $11{,}7810$ & $13{,}3237$ & $13{,}3518$ \\ \hline
$5$ & $14{,}9309$ & $14{,}9226$ & $16{,}4706$ & $16{,}4934$ \\ \hline
\end{tabular}
\end{table}
 
\begin{figure}[H]
\centering
\begin{tikzpicture}
\begin{axis}[bbeosc, width=.96\linewidth, height=6cm, xmin=0, xmax=12, ymin=-0.6, ymax=1.1, legend pos=north east, samples=300]
\bbezero{0}{12}
\addplot[black, very thick, smooth] coordinates {(0,1.0000) (0.1,0.9975) (0.2,0.9900) (0.3,0.9776) (0.4,0.9604) (0.5,0.9385) (0.6,0.9120) (0.7,0.8812) (0.8,0.8463) (0.9,0.8075) (1,0.7652) (1.1,0.7196) (1.2,0.6711) (1.3,0.6201) (1.4,0.5669) (1.5,0.5118) (1.6,0.4554) (1.7,0.3980) (1.8,0.3400) (1.9,0.2818) (2,0.2239) (2.1,0.1666) (2.2,0.1104) (2.3,0.0555) (2.4,0.0025) (2.5,-0.0484) (2.6,-0.0968) (2.7,-0.1424) (2.8,-0.1850) (2.9,-0.2243) (3,-0.2601) (3.1,-0.2921) (3.2,-0.3202) (3.3,-0.3443) (3.4,-0.3643) (3.5,-0.3801) (3.6,-0.3918) (3.7,-0.3992) (3.8,-0.4026) (3.9,-0.4018) (4,-0.3971) (4.1,-0.3887) (4.2,-0.3766) (4.3,-0.3610) (4.4,-0.3423) (4.5,-0.3205) (4.6,-0.2961) (4.7,-0.2693) (4.8,-0.2404) (4.9,-0.2097) (5,-0.1776) (5.1,-0.1443) (5.2,-0.1103) (5.3,-0.0758) (5.4,-0.0412) (5.5,-0.0068) (5.6,0.0270) (5.7,0.0599) (5.8,0.0917) (5.9,0.1220) (6,0.1506) (6.1,0.1773) (6.2,0.2017) (6.3,0.2238) (6.4,0.2433) (6.5,0.2601) (6.6,0.2740) (6.7,0.2851) (6.8,0.2931) (6.9,0.2981) (7,0.3001) (7.1,0.2991) (7.2,0.2951) (7.3,0.2882) (7.4,0.2786) (7.5,0.2663) (7.6,0.2516) (7.7,0.2346) (7.8,0.2154) (7.9,0.1944) (8,0.1717) (8.1,0.1475) (8.2,0.1222) (8.3,0.0960) (8.4,0.0692) (8.5,0.0419) (8.6,0.0146) (8.7,-0.0125) (8.8,-0.0392) (8.9,-0.0653) (9,-0.0903) (9.1,-0.1142) (9.2,-0.1367) (9.3,-0.1577) (9.4,-0.1768) (9.5,-0.1939) (9.6,-0.2090) (9.7,-0.2218) (9.8,-0.2323) (9.9,-0.2403) (10,-0.2459) (10.1,-0.2490) (10.2,-0.2496) (10.3,-0.2477) (10.4,-0.2434) (10.5,-0.2366) (10.6,-0.2276) (10.7,-0.2164) (10.8,-0.2032) (10.9,-0.1881) (11,-0.1712) (11.1,-0.1528) (11.2,-0.1330) (11.3,-0.1121) (11.4,-0.0902) (11.5,-0.0677) (11.6,-0.0446) (11.7,-0.0213) (11.8,0.0020) (11.9,0.0250) (12,0.0477) (12.1,0.0697) (12.2,0.0908) (12.3,0.1108) (12.4,0.1296) (12.5,0.1469) (12.6,0.1626) (12.7,0.1766) (12.8,0.1887) (12.9,0.1988) (13,0.2069) (13.1,0.2129) (13.2,0.2167) (13.3,0.2183) (13.4,0.2177) (13.5,0.2150) (13.6,0.2101) (13.7,0.2032) (13.8,0.1943) (13.9,0.1836) (14,0.1711) (14.1,0.1570) (14.2,0.1414) (14.3,0.1245) (14.4,0.1065) (14.5,0.0875) (14.6,0.0679) (14.7,0.0476) (14.8,0.0271) (14.9,0.0064) (15,-0.0142)};
\addplot[bbered, dashed, thick, domain=0.6:12] {sqrt(2/(pi*x))*cos(deg(x-pi/4))};
\legend{$J_0$, $\sqrt{2/(\pi x)}\cos(x-\pi/4)$}
\end{axis}
\end{tikzpicture}
\caption{$J_0$ y su aproximación asintótica: ya a partir de $x\approx3$ las curvas son casi indistinguibles.}
\label{fig:s5-asint}
\end{figure}
 
\begin{table}[H]
\centering
\small
\caption{Valores de las funciones de Bessel de orden $0$ y $1$.}
\label{tab:s5-valores}
\begin{tabular}{|c|r|r|r|r|}
\hline
\rowcolor{bbe!15}$x$ & $J_0(x)$ & $J_1(x)$ & $Y_0(x)$ & $Y_1(x)$ \\ \hline
$1$ & $0{,}7652$ & $0{,}4401$ & $0{,}0883$ & $-0{,}7812$ \\ \hline
$2$ & $0{,}2239$ & $0{,}5767$ & $0{,}5104$ & $-0{,}1070$ \\ \hline
$3$ & $-0{,}2601$ & $0{,}3391$ & $0{,}3769$ & $0{,}3247$ \\ \hline
$4$ & $-0{,}3971$ & $-0{,}0660$ & $-0{,}0169$ & $0{,}3979$ \\ \hline
$5$ & $-0{,}1776$ & $-0{,}3276$ & $-0{,}3085$ & $0{,}1479$ \\ \hline
$6$ & $0{,}1506$ & $-0{,}2767$ & $-0{,}2882$ & $-0{,}1750$ \\ \hline
$7$ & $0{,}3001$ & $-0{,}0047$ & $-0{,}0259$ & $-0{,}3027$ \\ \hline
$8$ & $0{,}1717$ & $0{,}2346$ & $0{,}2235$ & $-0{,}1581$ \\ \hline
$9$ & $-0{,}0903$ & $0{,}2453$ & $0{,}2499$ & $0{,}1043$ \\ \hline
$10$ & $-0{,}2459$ & $0{,}0435$ & $0{,}0557$ & $0{,}2490$ \\ \hline
\end{tabular}
\end{table}
 
\begin{example}\label{ex:s5-tambor}
\textbf{(Modos radiales de una membrana circular.)} Determine las soluciones acotadas de
\[ y''+\frac1x\,y'+\lambda^2y=0,\qquad 0<x<1,\qquad y(1)=0. \]
 
\noindent \textbf{Solución:} \quad Multiplicando por $x^2$ obtenemos $x^2y''+xy'+\lambda^2x^2y=0$, que es Bessel de orden $0$ con parámetro $\lambda$ (Ejemplo~\ref{ex:s5-parametrica}): $y=c_1J_0(\lambda x)+c_2Y_0(\lambda x)$. Como $Y_0$ no está acotada en $x=0$, necesitamos $c_2=0$. La condición $y(1)=0$ exige $J_0(\lambda)=0$, es decir, $\lambda=\lambda_k=j_{0,k}$: $\lambda_1\approx2{,}405$, $\lambda_2\approx5{,}520$, $\lambda_3\approx8{,}654$, \dots Las funciones propias son $y_k(x)=J_0(\lambda_kx)$ (Figura~\ref{fig:s5-modos}); la $k$-ésima tiene $k-1$ ceros en el interior de $(0,1)$ y es el perfil radial del $k$-ésimo modo de vibración de un tambor. \hfill $\square$
\end{example}
 
\begin{figure}[H]
\centering
\begin{tikzpicture}
\begin{axis}[bbeosc, width=.96\linewidth, height=6cm, xmin=0, xmax=1, ymin=-0.5, ymax=1.1, legend pos=north east, legend columns=3]
\bbezero{0}{1}
\addplot[bbeblue, very thick, smooth] coordinates {(0,1.0000) (0.01,0.9999) (0.02,0.9994) (0.03,0.9987) (0.04,0.9977) (0.05,0.9964) (0.06,0.9948) (0.07,0.9929) (0.08,0.9908) (0.09,0.9883) (0.1,0.9856) (0.11,0.9826) (0.12,0.9793) (0.13,0.9757) (0.14,0.9719) (0.15,0.9677) (0.16,0.9633) (0.17,0.9587) (0.18,0.9537) (0.19,0.9485) (0.2,0.9430) (0.21,0.9372) (0.22,0.9312) (0.23,0.9250) (0.24,0.9184) (0.25,0.9117) (0.26,0.9046) (0.27,0.8973) (0.28,0.8898) (0.29,0.8821) (0.3,0.8741) (0.31,0.8658) (0.32,0.8573) (0.33,0.8486) (0.34,0.8397) (0.35,0.8306) (0.36,0.8212) (0.37,0.8117) (0.38,0.8019) (0.39,0.7919) (0.4,0.7817) (0.41,0.7713) (0.42,0.7608) (0.43,0.7500) (0.44,0.7391) (0.45,0.7280) (0.46,0.7167) (0.47,0.7052) (0.48,0.6936) (0.49,0.6819) (0.5,0.6699) (0.51,0.6579) (0.52,0.6456) (0.53,0.6333) (0.54,0.6208) (0.55,0.6082) (0.56,0.5955) (0.57,0.5826) (0.58,0.5697) (0.59,0.5566) (0.6,0.5434) (0.61,0.5302) (0.62,0.5168) (0.63,0.5034) (0.64,0.4899) (0.65,0.4763) (0.66,0.4627) (0.67,0.4490) (0.68,0.4352) (0.69,0.4214) (0.7,0.4076) (0.71,0.3937) (0.72,0.3798) (0.73,0.3658) (0.74,0.3519) (0.75,0.3379) (0.76,0.3239) (0.77,0.3099) (0.78,0.2959) (0.79,0.2819) (0.8,0.2680) (0.81,0.2540) (0.82,0.2401) (0.83,0.2262) (0.84,0.2123) (0.85,0.1985) (0.86,0.1848) (0.87,0.1710) (0.88,0.1574) (0.89,0.1438) (0.9,0.1303) (0.91,0.1168) (0.92,0.1035) (0.93,0.0902) (0.94,0.0770) (0.95,0.0639) (0.96,0.0509) (0.97,0.0380) (0.98,0.0252) (0.99,0.0125) (1,0.0000)};
\addplot[bbered, very thick, smooth] coordinates {(0,1.0000) (0.01,0.9992) (0.02,0.9970) (0.03,0.9932) (0.04,0.9878) (0.05,0.9810) (0.06,0.9728) (0.07,0.9630) (0.08,0.9518) (0.09,0.9392) (0.1,0.9253) (0.11,0.9099) (0.12,0.8933) (0.13,0.8753) (0.14,0.8562) (0.15,0.8358) (0.16,0.8143) (0.17,0.7917) (0.18,0.7680) (0.19,0.7433) (0.2,0.7177) (0.21,0.6912) (0.22,0.6639) (0.23,0.6358) (0.24,0.6071) (0.25,0.5776) (0.26,0.5477) (0.27,0.5172) (0.28,0.4862) (0.29,0.4549) (0.3,0.4233) (0.31,0.3915) (0.32,0.3595) (0.33,0.3274) (0.34,0.2953) (0.35,0.2632) (0.36,0.2313) (0.37,0.1995) (0.38,0.1680) (0.39,0.1367) (0.4,0.1059) (0.41,0.0755) (0.42,0.0456) (0.43,0.0163) (0.44,-0.0124) (0.45,-0.0404) (0.46,-0.0677) (0.47,-0.0942) (0.48,-0.1198) (0.49,-0.1446) (0.5,-0.1684) (0.51,-0.1912) (0.52,-0.2131) (0.53,-0.2338) (0.54,-0.2535) (0.55,-0.2720) (0.56,-0.2894) (0.57,-0.3056) (0.58,-0.3206) (0.59,-0.3344) (0.6,-0.3469) (0.61,-0.3582) (0.62,-0.3682) (0.63,-0.3770) (0.64,-0.3844) (0.65,-0.3906) (0.66,-0.3955) (0.67,-0.3991) (0.68,-0.4015) (0.69,-0.4027) (0.7,-0.4025) (0.71,-0.4012) (0.72,-0.3987) (0.73,-0.3950) (0.74,-0.3902) (0.75,-0.3842) (0.76,-0.3772) (0.77,-0.3691) (0.78,-0.3600) (0.79,-0.3500) (0.8,-0.3390) (0.81,-0.3271) (0.82,-0.3143) (0.83,-0.3008) (0.84,-0.2865) (0.85,-0.2715) (0.86,-0.2559) (0.87,-0.2397) (0.88,-0.2229) (0.89,-0.2057) (0.9,-0.1880) (0.91,-0.1699) (0.92,-0.1516) (0.93,-0.1329) (0.94,-0.1141) (0.95,-0.0951) (0.96,-0.0761) (0.97,-0.0570) (0.98,-0.0379) (0.99,-0.0189) (1,0.0000)};
\addplot[bbegreen, very thick, smooth] coordinates {(0,1.0000) (0.01,0.9981) (0.02,0.9925) (0.03,0.9832) (0.04,0.9703) (0.05,0.9537) (0.06,0.9337) (0.07,0.9103) (0.08,0.8837) (0.09,0.8540) (0.1,0.8214) (0.11,0.7860) (0.12,0.7480) (0.13,0.7078) (0.14,0.6654) (0.15,0.6211) (0.16,0.5752) (0.17,0.5279) (0.18,0.4794) (0.19,0.4301) (0.2,0.3802) (0.21,0.3299) (0.22,0.2796) (0.23,0.2295) (0.24,0.1798) (0.25,0.1308) (0.26,0.0828) (0.27,0.0360) (0.28,-0.0094) (0.29,-0.0531) (0.3,-0.0950) (0.31,-0.1347) (0.32,-0.1723) (0.33,-0.2073) (0.34,-0.2399) (0.35,-0.2697) (0.36,-0.2966) (0.37,-0.3207) (0.38,-0.3417) (0.39,-0.3597) (0.4,-0.3745) (0.41,-0.3862) (0.42,-0.3948) (0.43,-0.4003) (0.44,-0.4026) (0.45,-0.4020) (0.46,-0.3984) (0.47,-0.3919) (0.48,-0.3826) (0.49,-0.3707) (0.5,-0.3563) (0.51,-0.3395) (0.52,-0.3206) (0.53,-0.2996) (0.54,-0.2768) (0.55,-0.2524) (0.56,-0.2265) (0.57,-0.1994) (0.58,-0.1713) (0.59,-0.1424) (0.6,-0.1130) (0.61,-0.0831) (0.62,-0.0532) (0.63,-0.0233) (0.64,0.0062) (0.65,0.0353) (0.66,0.0636) (0.67,0.0911) (0.68,0.1174) (0.69,0.1426) (0.7,0.1663) (0.71,0.1884) (0.72,0.2088) (0.73,0.2274) (0.74,0.2440) (0.75,0.2586) (0.76,0.2711) (0.77,0.2814) (0.78,0.2894) (0.79,0.2953) (0.8,0.2988) (0.81,0.3001) (0.82,0.2991) (0.83,0.2960) (0.84,0.2906) (0.85,0.2832) (0.86,0.2737) (0.87,0.2624) (0.88,0.2491) (0.89,0.2342) (0.9,0.2177) (0.91,0.1998) (0.92,0.1806) (0.93,0.1602) (0.94,0.1389) (0.95,0.1168) (0.96,0.0940) (0.97,0.0707) (0.98,0.0472) (0.99,0.0236) (1,0.0000)};
\legend{$J_0(\lambda_1x)$, $J_0(\lambda_2x)$, $J_0(\lambda_3x)$}
\end{axis}
\end{tikzpicture}
\caption{Ejemplo~\ref{ex:s5-tambor}: los tres primeros modos radiales, que se anulan en el borde $x=1$.}
\label{fig:s5-modos}
\end{figure}
 
\begin{exercise}[Para llevar]
\begin{enumerate}
\item Exprese $J_4(x)$ en términos de $J_0(x)$ y $J_1(x)$.
\item Encuentre la solución general de $x^2y''+xy'+(4x^2-9)y=0$, $x>0$, y la solución acotada en $x=0$.
\item Demuestre que $J_{1/2}'(x)=\dfrac12\bigl[J_{-1/2}(x)-J_{3/2}(x)\bigr]$ coincide con la derivada de $\sqrt{2/(\pi x)}\sin x$ (Propiedad 4 y Ejemplo~\ref{ex:s5-medio}).
\end{enumerate}
\emph{Respuestas:} (1) $J_4=\dfrac{6}{x}J_3-J_2=\Bigl(\dfrac{48}{x^3}-\dfrac8x\Bigr)J_1(x)+\Bigl(1-\dfrac{24}{x^2}\Bigr)J_0(x)$. \; (2) Con $z=2x$ es Bessel de orden $3$: $y=c_1J_3(2x)+c_2Y_3(2x)$; la solución acotada es $c_1J_3(2x)$. \; (3) La Propiedad 4 es la identidad general $J_p'=\frac12(J_{p-1}-J_{p+1})$; con $p=\tfrac12$ y $J_{3/2}=\frac1x J_{1/2}-J_{-1/2}$ (Propiedad 3) se obtiene $J_{1/2}'=J_{-1/2}-\frac{1}{2x}J_{1/2}$, que es exactamente la derivada de $\sqrt{2/(\pi x)}\sin x$.
\end{exercise}
